% Figura: línea temporal (no a escala) del electromagnetismo hasta 1905.
\begin{figure}[ht]
\centering
\begin{tikzpicture}[x=1.52cm, y=1cm]
  \draw[line width=1.1pt, azulacero, -{Stealth[length=6pt]}] (-0.4,0) -- (9.5,0);
  \tikzset{ev/.style={font=\scriptsize, align=center, text width=2.55cm, text=black!85},
           anio/.style={font=\scriptsize\bfseries, text=azulprof}}
  \foreach \i/\an/\txt/\pos in {
      0/{s.\,VI a.\,C.}/{Tales de Mileto:\\el ámbar frotado}/1,
      1/{1600}/{Gilbert,\\\emph{De Magnete}}/-1,
      2/{1771}/{Cavendish:\\ley del cuadrado}/1,
      3/{1785}/{Coulomb:\\balanza de torsión}/-1,
      4/{1800}/{Volta:\\la pila}/1,
      5/{1820}/{Oersted y Ampère}/-1,
      6/{1831}/{Faraday:\\inducción y campo}/1,
      7/{1873}/{Maxwell:\\\emph{Tratado}}/-1,
      8/{1888}/{Hertz:\\ondas e.\,m.}/1,
      9/{1900--05}/{Planck,\\Einstein}/-1}{
    \fill[azulprof] (\i,0) circle (2.2pt);
    \ifnum\pos=1
      \draw[grispizarra!60, line width=0.4pt] (\i,0.1) -- (\i,0.55);
      \node[anio, anchor=south] at (\i,0.52) {\an};
      \node[ev, anchor=south] at (\i,0.95) {\txt};
    \else
      \draw[grispizarra!60, line width=0.4pt] (\i,-0.1) -- (\i,-0.55);
      \node[anio, anchor=north] at (\i,-0.52) {\an};
      \node[ev, anchor=north] at (\i,-0.95) {\txt};
    \fi}
  \node[font=\tiny, text=grispizarra, anchor=north east] at (9.45,-0.12) {no a escala};
\end{tikzpicture}
\caption{Hitos del electromagnetismo, de la Antigüedad al cambio de siglo. El eje
no está a escala: casi todo ocurre entre 1771 y 1888.}
\label{fig:historia}
\end{figure}
